% C++17 language and library lookup; examples use the contest template.
\section{C++17 language reference}
Examples assume \texttt{using namespace std;} and
\texttt{using ll = long long;}. \texttt{<bits/stdc++.h>} is a GCC convenience
header; it is not part of ISO C++. Prefer Python for exact huge integers.

\subsection{Input, output and terminal debugging}\label{cpp:io}
\begin{lstlisting}[style=cpp]
ios::sync_with_stdio(false);
cin.tie(nullptr);
int n; string word;
cin >> n >> word;                 // whitespace-separated
vector<int> a(n);
for (int& x : a) cin >> x;
cout << n << ' ' << word << '\n';
cout << fixed << setprecision(10) << 1.0 / 3 << '\n';
int x;
while (cin >> x) { /* integers until EOF */ }
\end{lstlisting}
\hint{Use \texttt{'\textbackslash n'} for ordinary output;
\texttt{endl} also flushes. Interactive output must flush:
\texttt{cout << "? " << x << endl;}. After disabling synchronization,
avoid mixing \texttt{cin/cout} with \texttt{scanf/printf}.}
\begin{lstlisting}[style=cpp]
int n; cin >> n;
cin.ignore(numeric_limits<streamsize>::max(), '\n');
string line;
getline(cin, line);               // next line, spaces retained
// getline(cin >> ws, line);       // skips ALL leading whitespace
cerr << "n = " << n << '\n';      // stderr: stays in terminal
dbg(n);                          // template: only with -DLOCAL
\end{lstlisting}
\hint{\texttt{>>} leaves the newline behind. \texttt{getline} removes its
delimiter but preserves other spaces; \texttt{>> ws} also skips blank lines.
Redirecting stdout with \texttt{> out.txt} leaves stderr visible.
Raw \texttt{cerr} always runs; \texttt{dbg} is conditional on \texttt{LOCAL}.}

\subsection{Declarations, loops and references}\label{cpp:syntax}
\begin{lstlisting}[style=cpp]
using ll = long long;
using pii = pair<int, int>;
vector<int> a = {3, 1, 4};
const int MOD = 1000000007;
auto n = a.size();                // unsigned size type
for (int i = 0; i < int(a.size()); ++i) { /* a[i] */ }
for (int x : a) { /* copy of each value */ }
for (int& x : a) ++x;             // change original elements
for (const auto& x : a) { /* read without copying */ }
auto copy = a;                    // independent vector copy
auto& alias = a;                  // another name for a
swap(a[0], a[1]);
\end{lstlisting}
\hint{\texttt{const T\&} passes a read-only reference; \texttt{T\&} allows
mutation. Passing \texttt{vector<T>} by value copies it. Local scalar
variables are uninitialized until assigned. Indices run from 0 to size$-1$;
\texttt{v[i]} does not check bounds, \texttt{v.at(i)} does.}
\begin{lstlisting}[style=cpp]
ll square(ll x) { return x * x; }  // define outside main
void add_one(vector<int>& a) {
    for (int& x : a) ++x;
}
int offset = 5;
auto shifted = [offset](int x) { return x + offset; };
auto change = [&](int x) { offset += x; };  // capture by ref
auto twice = [](auto x) { return x + x; }; // generic lambda
\end{lstlisting}
\hint{\texttt{[x]} copies the current value of \texttt{x};
\texttt{[\&x]} refers to the original. A reference capture must not outlive
its variable. Use \texttt{struct State \{ int x, y; \};} for named fields.}

\subsection{Vectors, grids, pairs and tuples}\label{cpp:containers}
\begin{lstlisting}[style=cpp]
vector<int> v(5, -1);             // five copies of -1
v.push_back(7); v.pop_back();
v.resize(8, 0);                  // change size; fill new items
v.reserve(100);                  // capacity only; no new items
v.assign(3, 9);                  // replace with [9,9,9]
v.insert(v.begin() + 1, 4);       // O(n), insert before index 1
v.erase(v.begin() + 1);          // O(n), remove index 1
int first = v.front(), last = v.back(); // require nonempty
vector<vector<int>> grid(3, vector<int>(4, 0));
pair<int, int> p = {2, 7};
auto [x, y] = p;                 // C++17 structured binding
tuple<int, int, string> t = {1, 2, "a"};
auto [u, w, name] = t;
get<0>(t) = 9;                   // index is compile-time
tie(x, y) = p;                   // assign existing variables
\end{lstlisting}
\hint{Each grid row is independent. \texttt{auto [x,y]} copies fields;
\texttt{auto\& [x,y]} refers to them. Pairs, tuples, vectors and strings
compare lexicographically. Reallocating a vector invalidates its references
and iterators; erasing invalidates positions at or after the erased element.}

\subsection{Strings and characters}\label{cpp:strings}
\begin{lstlisting}[style=cpp]
string s = "banana";
s += '!'; s.pop_back();
s.substr(1, 3);                  // "ana": start, length
auto pos = s.find("na");         // 2; string::npos if absent
if (pos != string::npos) s.replace(pos, 2, "X");
s.insert(0, "pre"); s.erase(0, 3);
reverse(s.begin(), s.end());
string digits = to_string(123);
ll value = stoll("-123");        // decimal by default
int d = '7' - '0';               // digit character -> 7
char c = 'A';
c = char(tolower((unsigned char)c));
bool digit = isdigit((unsigned char)c);
\end{lstlisting}
\hint{C++ strings are mutable byte sequences; single quotes create a
character, double quotes a string. \texttt{substr} takes a length, not an
end index. Check \texttt{npos} before converting a found position to
\texttt{int}. \texttt{stoi/stoll} throw on invalid or out-of-range input.}

\subsection{Ordered sets, maps and multisets}
\begin{lstlisting}[style=cpp]
set<int> s = {2, 5, 9};           // sorted, unique
s.insert(4); s.erase(9);
bool found = s.find(5) != s.end(); // C++17: no contains()
auto it = s.lower_bound(3);       // first >= 3, or end()
if (it != s.end()) cout << *it << '\n';
if (!s.empty()) cout << *s.begin() << ' ' << *s.rbegin();
multiset<int> ms = {2, 2, 7};
auto one = ms.find(2);
if (one != ms.end()) ms.erase(one); // erase ONE occurrence
ms.erase(2);                     // erase ALL remaining 2s
map<string, int> counts;
counts["apple"]++;               // missing key inserts value 0
auto hit = counts.find("pear");   // lookup without insertion
if (hit != counts.end()) cout << hit->second;
for (const auto& [key, val] : counts) { /* sorted keys */ }
map<pair<int, int>, int> at;      // pair keys work directly
\end{lstlisting}
\hint{Ordered lookup/insert costs $O(\log n)$. Iterators are not indices:
\texttt{it - s.begin()} is invalid, and \texttt{distance} is $O(n)$.
Use \texttt{s.lower\_bound(x)}, not the generic algorithm on set iterators.
Never dereference \texttt{end()} or erase a failed \texttt{find}.}
\begin{lstlisting}[style=cpp]
unordered_map<string, int> freq; // average O(1) lookup
freq.reserve(10000);
freq["a"]++;
bool has = freq.find("a") != freq.end();
unordered_set<int> seen;
seen.insert(4); seen.erase(4);
\end{lstlisting}
\hint{Hash containers have no sorted order; worst-case operations can be
$O(n)$. Rehashing invalidates iterators. \texttt{pair} and \texttt{vector}
have no standard C++17 hash; use an ordered map or supply a custom hash.}

\subsection{Queues, stacks, deques and heaps}
\begin{lstlisting}[style=cpp]
deque<int> d = {1, 2};
d.push_front(0); d.push_back(3);
d.pop_front(); d.pop_back();      // O(1) ends; d[i] is O(1)
queue<int> q;
q.push(4); int a = q.front(); q.pop();
stack<int> st;
st.push(5); int b = st.top(); st.pop();
priority_queue<int> maxq;
priority_queue<int, vector<int>, greater<int>> minq;
maxq.push(7); maxq.push(2);
int largest = maxq.top(); maxq.pop();
using Item = pair<ll, int>;       // distance, vertex
priority_queue<Item, vector<Item>, greater<Item>> heap;
heap.push({12, 3});
auto [dist, vertex] = heap.top(); heap.pop();
\end{lstlisting}
\hint{All \texttt{pop()} operations return nothing: read the value first.
Check \texttt{empty()} before \texttt{front/top/pop}. Heap push/pop costs
$O(\log n)$; top is $O(1)$. Pair heaps break ties by the second field.}

\subsection{Sorting and binary search}\label{cpp:order}
\begin{lstlisting}[style=cpp]
vector<int> v = {5, 1, 4, 1};
sort(v.begin(), v.end());         // [1,1,4,5]
sort(v.begin(), v.end(), greater<int>()); // descending
stable_sort(v.begin(), v.end());  // equal keys keep input order
vector<pair<int, int>> a = {{1, 9}, {2, 3}, {3, 3}};
sort(a.begin(), a.end(), [](const auto& x, const auto& y) {
    if (x.second != y.second) return x.second < y.second;
    return x.first > y.first;    // tie: first field descending
});
\end{lstlisting}
\hint{A comparator must give a strict weak ordering: \texttt{cmp(x,x)}
must be false; use \texttt{<}, not \texttt{<=}. Do not subtract keys to
compare them: subtraction can overflow.}
\begin{lstlisting}[style=cpp]
vector<int> a = {1, 2, 2, 5};      // sorted ascending
int x = 2;
auto lo = lower_bound(a.begin(), a.end(), x); // first >= x
auto hi = upper_bound(a.begin(), a.end(), x); // first > x
int first = int(lo - a.begin());  // 1; size() if none
int count = int(hi - lo);         // 2 equal values
bool exists = binary_search(a.begin(), a.end(), x);
int k = 2;                       // zero-based, 0 <= k < size
nth_element(a.begin(), a.begin() + k, a.end());
int kth = a[k];                  // kth smallest; not fully sorted
\end{lstlisting}
\hint{Binary search needs a range sorted using the same comparison.
Iterator ranges are half-open: \texttt{[begin,end)}. Vector binary search
costs $O(\log n)$; insertion or erasure in its middle costs $O(n)$.}

\subsection{Algorithm utilities}
\begin{lstlisting}[style=cpp]
vector<int> v = {3, 1, 3, 4};
ll sum = accumulate(v.begin(), v.end(), 0LL);
int mx = *max_element(v.begin(), v.end()); // nonempty
int threes = count(v.begin(), v.end(), 3);
bool has = find(v.begin(), v.end(), 1) != v.end();
fill(v.begin(), v.end(), 7);
iota(v.begin(), v.end(), 0);      // 0,1,2,3
reverse(v.begin(), v.end());
v.erase(remove(v.begin(), v.end(), 2), v.end()); // drop all 2s
sort(v.begin(), v.end());
v.erase(unique(v.begin(), v.end()), v.end());   // deduplicate
ll g = gcd(12LL, 18LL), l = lcm(4LL, 6LL);
vector<int> p = {0, 1, 2};        // start sorted
do { /* use p */ } while (next_permutation(p.begin(), p.end()));
\end{lstlisting}
\hint{\texttt{remove} and \texttt{unique} return a new logical end;
\texttt{erase} actually shrinks the vector. \texttt{unique} only merges
adjacent equal values. The initial value of \texttt{accumulate} selects
its arithmetic type: \texttt{0LL} for \texttt{long long}.}

\subsection{Integer arithmetic, limits and modulo}\label{cpp:numbers}
\begin{lstlisting}[style=cpp]
int x = 1000000000, y = 1000000000;
ll product = 1LL * x * y;         // widen BEFORE multiplication
ll big = numeric_limits<ll>::max();
const ll INF = (1LL << 62);       // chosen sentinel, not infinity
int a = -7, b = 3;
int q = a / b, rem = a % b;       // -2, -1: division truncates
// For integer a, positive b:
int floor_div = a / b - (a % b < 0);
int ceil_div = a / b + (a % b > 0);
ll m = 1000000007, value = -12;   // m > 0
ll r = value % m; if (r < 0) r += m; // normalize to [0,m)
__int128 wide = (__int128)product * product; // GCC extension
ll reduced = (ll)(wide % m);
\end{lstlisting}
\hint{Signed overflow is undefined behavior, including intermediate
expressions. Assigning an overflowing \texttt{int} product to \texttt{ll}
does not repair it. Integer division truncates even when assigned to
\texttt{double}: use \texttt{double(a)/b}. \texttt{pow} uses floating point;
it is not exact integer exponentiation or modular exponentiation.
C++17 has no standard arbitrary-precision integer type.}

\subsection{Bitmasks and bit counts}
\begin{lstlisting}[style=cpp]
unsigned long long mask = 0;
int i = 5;                       // require 0 <= i < 64
mask |= 1ULL << i;               // set
mask &= ~(1ULL << i);            // clear
mask ^= 1ULL << i;               // toggle
bool set = (mask & (1ULL << i)) != 0;
int ones = __builtin_popcountll(mask); // GCC, 64-bit operand
if (mask) {
    int low = __builtin_ctzll(mask);  // lowest set bit index
    int high = 63 - __builtin_clzll(mask); // highest set bit
}
bitset<128> bits;
bits.set(7); bits.flip(2); bits.reset(7);
auto count = bits.count();
\end{lstlisting}
\hint{\texttt{ctz/clz} require a nonzero operand. Use unsigned operands
for bit operations; shifting by a negative amount or by the type's width
is invalid. The builtins above are GCC extensions; C++20
\texttt{popcount/countr\_zero} are outside the C++17 baseline.}
